%------------------------------------------------------------------------------%

\usepackage{integracao_e_series}
\usepackage{cancel}

\title{Séries Numéricas -- Definição e Exemplos}

%------------------------------------------------------------------------------%
\begin{document}

\addtitle

%------------------------------------------------------------------------------%
\section{Definição}

%------------------------------------------------------------------------------%
\begin{frame}{Séries}
    Uma \emph{Série Infinita} é a soma dos termos de uma \emph{Sequência Infinita} $(a_n)$

    \pause
    \vspace{\baselineskip}
    Se existir, essa \emph{soma} pode ser escrita como
    \[
      \sum_{n=1}^\infty a_n
      \;=\; a_1 + a_2 + \cdots + a_n + \cdots
      \;=\; S
    \]

    \pause
    Cada valor $a_n$ é chamado de \emph{Termo da Série}

		\only<beamer>{
    \pause
    \vspace{\baselineskip}
    \begin{center}
        \textbf{\color{red}Como saber se a soma existe?}
    \end{center}}
\end{frame}

%------------------------------------------------------------------------------%
\begin{frame}{Sequência de Somas Parciais}
    Soma dos primeiros $n$ \emph{termos} de uma sequência $(a_k)$
    \[
      S_{\memph{n}}
      = \; a_1 + a_2 + \cdots + a_{\memph{n}} \;
      = \sum_{k=1}^{\memph{n}} a_k
    \]

    \pause
    Temos agora uma nova sequência: \emph{Somas Parciais} $(S_n)$

    \pause
    \vspace{\baselineskip}
    Podemos usar as técnicas definidas para sequências

\end{frame}

%------------------------------------------------------------------------------%
\begin{frame}{Convergência de uma Série}
    A \emph{soma da série} é $S$
    \[
      \sum_{n=1}^\infty a_n = S
    \]
    quando a sequência de \emph{somas parciais} $(S_n)$ converge para $S$
    \pause
    \[
      \lim_{n\to\infty} S_n
      \;=\; \lim_{n\to\infty} \sum_{k=1}^n a_k
      \;=\; S
    \]
    \pause
    Se sequência de somas parciais diverge a série \emph{diverge}
\end{frame}

%------------------------------------------------------------------------------%
\begin{frame}{Exemplo}
    Somar os termos da sequência
    \[ a_k = \frac{1}{2^k} \]

    \pause
    \vspace{.5\baselineskip}
    \[
    	a_1 = \frac{1}{2}  \qquad\pause
    	a_2 = \frac{1}{4}  \qquad\pause
    	a_3 = \frac{1}{8}  \qquad\pause
      a_4 = \frac{1}{16} \qquad\pause
      a_5 = \frac{1}{32} \qquad\pause\cdots
    \]
\end{frame}

%------------------------------------------------------------------------------%
\begin{frame}{Somas Parciais}
    \begin{align*}
      \onslide<+->{a_1 & = \frac{1}{2}   & }
      \onslide<+->{S_1 & = \frac{1}{2}                             &  & }
      \onslide<+->{\hspace{-17mm}  = 1 - \frac{1}{2}   \\[2mm]}
    	\onslide<+->{a_2 & = \frac{1}{4}   & }
      \onslide<+->{S_2 & = \frac{1}{2} + \frac{1}{4}               &  & }
      \onslide<+->{\hspace{-17mm}  = 1 - \frac{1}{4}   \\[2mm]}
    	\onslide<+->{a_3 & = \frac{1}{8}   & }
      \onslide<+->{S_3 & = \frac{1}{2} + \frac{1}{4} + \frac{1}{8} &  & }
      \onslide<+->{\hspace{-17mm}  = 1 - \frac{1}{8}   \\[2mm]}
    	\onslide<+->{a_4 & = \frac{1}{16}  & }
      \onslide<+->{S_4 & = \frac{1}{2} + \frac{1}{4} + \frac{1}{8} + \frac{1}{16} &  & }
      \onslide<+->{\hspace{-17mm}  = 1 - \frac{1}{16}   \\[7mm]}
      \onslide<+->{a_n & = \frac{1}{2^n} & }
      \onslide<+->{S_n & = \sum_{k=1}^{n} \frac{1}{2^k}        &  & }
      \onslide<+->{\hspace{-17mm}  = \emph{1 - \frac{1}{2^n} \quad \text{(Conjectura)} }}
    \end{align*}
\end{frame}

%------------------------------------------------------------------------------%
\begin{frame}{Demonstração dor Indução Finita}
  Demonstração dor Indução Finita

  \pause
  \vspace{\baselineskip}
	Para mostrar que uma afirmação é verdadeira para todo $n\in\N$

  devemos provar que
  \pause
	\begin{enumerate}[<+->]
		\item ela é verdadeira para $n=1$
		\item se ela for verdadeira para $n$ então também será para $n+1$
	\end{enumerate}
\end{frame}

%------------------------------------------------------------------------------%
\begin{frame}{Demonstração da Fórmula das Somas Parciais}
\begin{columns}
\column{0.5\linewidth}
  Assumimos por \emph{indução finita} que
  \quad
  \[
    S_{n-1} = 1 - \frac{1}{2^{n-1}}
  \]
  para um \emph{$n$ dado qualquer}

  \pause
  \vspace{\baselineskip}
  Queremos mostrar que
  \[
    S_n = 1 - \frac{1}{2^n}
  \]

  \pause
  \column{0.5\linewidth}
  \begin{align*}
    \onslide<+->{S_n & = \frac{1}{2}
                       + \frac{1}{4}
                       + \frac{1}{8}
                       + \cdots
                       + \frac{1}{2^{n-1}}
                       + \frac{1}{2^n}} \\[2mm]
      \onslide<+->{& = \left(\frac{1}{2}
                     + \frac{1}{4}
                     + \frac{1}{8}
                     + \cdots
                     + \frac{1}{2^{n-1}}\right)
                     + \frac{1}{2^n}} \\[2mm]
    \onslide<+->{& = {\color<5->{blue}S_{n-1}} + \frac{1}{2^n}\\[2mm]}
    \onslide<+->{& = {\color<5->{blue}1 - \frac{1}{2^{n-1}}} + \frac{1}{2^n}\\[2mm]}
    \onslide<+->{& = 1 - \frac{2}{2^n} + \frac{1}{2^n}\\[2mm]}
    \onslide<+->{& = 1 - \frac{1}{2^n}}
  \end{align*}
\end{columns}
\end{frame}

%------------------------------------------------------------------------------%
\begin{frame}{Calculando o Limite das Somas Parciais}
    \begin{align*}
    	\onslide<+->{\lim_{n\to\infty} S_n}
      \onslide<+->{& = \lim \left(1 - \frac{1}{2^n}\right)     \\[2mm]}
    	\onslide<+->{& = \lim 1 - \lim\left(\frac{1}{2}\right)^{\!\!n} \\[2mm]}
    	\onslide<+->{& = 1 - 0 \\[2mm]}
      \onslide<+->{& = 1}
    \end{align*}

    \onslide<+->
    {
      Portanto a série converge e podemos escrever
      \[
        \sum_{n=1}^{\infty} \frac{1}{\;2^n\;} = 2
      \]
    }

\end{frame}

%------------------------------------------------------------------------------%
%------------------------------------------------------------------------------%
\section{Séries Geométricas}
%------------------------------------------------------------------------------%
%------------------------------------------------------------------------------%

%------------------------------------------------------------------------------%
\begin{frame}{Séries Geométricas}
	Uma \emph{Série Geométrica} é uma série com a forma
	\[
    \alpha + \alpha r + \alpha r^2 + \alpha r^3 + \cdots + \alpha r^{n-1} + \cdots
    = \sum_{n=1}^\infty \alpha r^{n-1}
  \]

  \pause\qquad\(\alpha\in\R\)          \tabto{33mm} número fixo não nulo, $\alpha\neq 0$

	\vspace{0.5\baselineskip}
  \pause\qquad\(r\in\R\)               \tabto{33mm} \emph{razão} da série

  \vspace{0.5\baselineskip}
  \pause\qquad\(a_n = \alpha r^{n-1}\) \tabto{33mm} termo geral

	\pause\vspace{\baselineskip}
	Especial pois sabemos quando converge e qual sua soma
\end{frame}

%------------------------------------------------------------------------------%
\begin{frame}{Convergência da Série Geométrica -- Caso 1}
	Se $r=1$
  \pause
  temos
	\[
    S_n
    = \sum_{k=1}^{n} \alpha \,1^{k-1}
    \pause
    = \sum_{k=1}^{n} \alpha
    \pause
    = \underbrace{\alpha + \cdots + \alpha}_{n \text{ vezes}}
    \pause
    = n\alpha
  \]

	\pause
	\vspace{\baselineskip}
	Portanto $S_n$ diverge quando $n\to\infty$
\end{frame}

%------------------------------------------------------------------------------%
\begin{frame}{Convergência da Série Geométrica -- Caso 2}
	Se $r\neq 1$ temos
	\begin{align*}
		\onslide<+->{S_n      & = \alpha  +    \alpha r + \alpha r^2 + \cdots + \alpha r^{n-1}                \\[2mm]}
		\onslide<+->{rS_n     & = \hspace{7mm} \alpha r + \alpha r^2 + \cdots + \alpha r^{n-1} + \alpha r^{n} \\[2mm]}
		\onslide<+->{S_n-rS_n & = \alpha - \alpha r^{n} \\[2mm]}
		\onslide<+->{S_n(1-r) & = \alpha - \alpha r^{n}        }
	\end{align*}
	\onslide<+->{
    \[
      S_n = \frac{\alpha (1-r^n)}{1-r} \qquad\qquad (r\neq 1)
    \]
    }
\end{frame}

%------------------------------------------------------------------------------%
\begin{frame}{Calculando o Limite}
	Com $r\neq 1$, temos \(S_n = \frac{\alpha(1-r^n)}{1-r}\),
  \pause
  portanto
  \[
    \lim S_n
    \pause
    = \lim \frac{\alpha(1-r^n)}{1-r}
    \pause
    = \frac{\alpha}{1-r} -\frac{\alpha}{1-r} \lim r^n
  \]
  \pause
  \vspace{-\baselineskip}
	\begin{description}
		\item[$r=-1$] \tabto{10mm}
		\onslide<6->{$r^n=(-1)^n\;$ diverge quando $n\to\infty$ portanto $S_n$ diverge} \\[7mm]
		\item[$\abs{r}>1$] \tabto{10mm}
		\onslide<7->{$\abs{r^n}\to\infty$ quando $n\to\infty$ portanto $S_n$ diverge} \\[7mm]
		\item[$\abs{r}<1$] \tabto{10mm}
		\onslide<8->{$r^n\to 0$ quando $n\to\infty$ portanto \( S_n \to \frac{\alpha}{1-r} \)}
	\end{description}
\end{frame}

%------------------------------------------------------------------------------%
\begin{frame}{Séries Geométricas}
	Uma série geométrica
	\[
    \sum_{n=1}^\infty \alpha r^{n-1}
  \]
	é convergente se $\abs{r}<1$ com soma
	\[
    \sum_{n=1}^\infty \alpha r^{n-1} = \frac{\alpha }{1-r}
  \]
	e divergente se $\abs{r}\geq 1$
\end{frame}

%------------------------------------------------------------------------------%
%------------------------------------------------------------------------------%
\section{Séries Telescópicas}
%------------------------------------------------------------------------------%
%------------------------------------------------------------------------------%

%------------------------------------------------------------------------------%
\begin{frame}{Séries Telescópicas}
	Raramente conseguimos fórmulas exatas para a soma de séries

	\pause
	\vspace{\baselineskip}
	\emph{Séries Telescópicas} são um caso especial onde quase todos os termos se anulam
\end{frame}

%------------------------------------------------------------------------------%
\begin{frame}{Exemplo}
	Calcular a soma da série
	\[\sum_{n=1}^{\infty} \left( \frac{1}{n} - \frac{1}{n+1} \right) \]
\end{frame}

%------------------------------------------------------------------------------%
\begin{frame}{Somas Parciais}
	\begin{align*}
		\onslide<+->{S_n & = \sum_{k=1}^{n} \left( \frac{1}{k} - \frac{1}{k+1} \right)
			\\[3mm]}
		\onslide<+->{& = \left( \frac{1}{1} - \frac{1}{1+1} \right)
			+ \left( \frac{1}{2} - \frac{1}{2+1} \right)
			+ \cdots
			+ \left( \frac{1}{n-1} - \frac{1}{n} \right)
			+ \left( \frac{1}{n} - \frac{1}{n+1} \right)
			\\[3mm]}
		\onslide<+->{&
      = \frac{1}{1} - {\color<4->{blue}\frac{1}{2}}
			+ {\color<4->{blue}\frac{1}{2}} - {\color<5->{red}\frac{1}{3}}
      + {\color<5->{red}\frac{1}{3}}  - {\color<6->{gray}\frac{1}{4}}
			+ \cdots
      + {\color<6->{gray}\frac{1}{n-2}}  - {\color<7->{green}\frac{1}{n-1}}
			+ {\color<7->{green}\frac{1}{n-1}} - {\color<8->{magenta}\frac{1}{n}}
			+ {\color<8->{magenta}\frac{1}{n}} - \frac{1}{n+1}
			\\[3mm]}
		\onslide<9->{& = 1 - \frac{1}{n+1}}
	\end{align*}
\end{frame}

%------------------------------------------------------------------------------%
\begin{frame}{Calculando o Limite}
	\begin{align*}
		\onslide<+->{\sum_{n=1}^{\infty} \left( \frac{1}{n} - \frac{1}{n+1} \right)
                  & = \lim_{n\to\infty} S_n \\[3mm]}
		\onslide<+->{ & = \lim_{n\to\infty} \left( 1 - \frac{1}{n+1} \right)  \\[3mm]}
		\onslide<+->{ & = 1}
	\end{align*}
\end{frame}

%------------------------------------------------------------------------------%
\begin{frame}{Pode Não Ser Óbvio}
	Calcular a soma da série
	\[\sum_{n=1}^{\infty} \frac{1}{\;n(n+1)\;} \]
	\pause
	Usando frações parciais
	\[ \frac{1}{n(n+1)} = \frac{1}{n} - \frac{1}{n+1} \]
	\pause
	Temos
	\[
	\sum_{n=1}^{\infty} \frac{1}{\;n(n+1)\;}
	\;=\; \sum_{n=1}^{\infty} \left( \frac{1}{n} - \frac{1}{n+1} \right)
	\;=\; 1
	\]
\end{frame}

%------------------------------------------------------------------------------%
%------------------------------------------------------------------------------%
\section{Testes da Divergência}
%------------------------------------------------------------------------------%
%------------------------------------------------------------------------------%

%------------------------------------------------------------------------------%
\begin{frame}{Testes de Convergência}
	\begin{itemize}[<+->]
		\item Em geral não temos uma fórmula fechada para $S_n$ \\[2mm]
		\item Queremos saber se a série converge \\[2mm]
		\item Abrimos mão de calcular sua soma exata \\[2mm]
		\item Podemos calcular aproximações computacionalmente
	\end{itemize}
\end{frame}

%------------------------------------------------------------------------------%
\begin{frame}{Termo Geral de uma Serie Convergente}
	Como somamos infintos termos se eles não forem ``pequenos'' a série não converge

	\pause
	\vspace{\baselineskip}
	Para que \(\;\sum_{n=1}^\infty a_n\;\) seja convergente $a_n$ precisa ir para zero
\end{frame}

%------------------------------------------------------------------------------%
\begin{frame}{Limite do Termo Geral}

	Se a série \(\sum_{n=1}^\infty a_n\) é convergente, então $a_n\to 0$

	\pause
	\vspace{\baselineskip}
	Como a série converge
	\[ \lim S_n  = S \]

	\pause
	Calculando
	\begin{align*}
		\onslide<+->{\lim a_n & = \lim(S_n-S_{n-1})       \\}
		\onslide<+->{         & = \lim S_n - \lim S_{n-1} \\}
		\onslide<+->{         & = S - S = 0}
	\end{align*}
\end{frame}

%------------------------------------------------------------------------------%
\begin{frame}{Teste da Divergência}
	Se a sequência $a_n$ não converge para zero, então a série $\sum a_n$ diverge
\end{frame}

%------------------------------------------------------------------------------%
\section{Lista Mínima}

%------------------------------------------------------------------------------%
\begin{frame}{Lista Mínima}
  Estudar as Seções 6.1 , 6.2 e 6.3 da Apostila

  \vfill
  Exercícios:

  \quad Seção 6.1: 1, 2, 4a-c, 6a-b, 8a-c, 10, 11a-b

  \quad Seção 6.2: 1

  \quad Seção 6.3: 3a-d

  \vfill
  Atenção: A prova é baseada no livro, não nas apresentações
\end{frame}


%------------------------------------------------------------------------------%
\end{document}
%------------------------------------------------------------------------------%
